\def\ChapterCode{FORTBILDUNG}
	\chapter
[\CO{[FORTBILDUNG]} Fehlermanagement]
{Fehlermanagement}
\label{chp:ADVANCED-Fehlermanagement}

\ChapterIntro
{%

}
{%
\begin{description}
  \item[Maintainer:] \ldots
  \item[Autoren:] Diverse
  \item[Reviewer:] Standard-Reviewprozess
  \item[Version:] {\VarDocVersion} ({\VarDocVersionDescription})
  \item[SHA1:] {\DocGitCommit}
\end{description}

{\TxTeilkapitelAASS}}



\clearpage % TEMP temp clearpage



[Warning: Draw object ignored][Warning: Draw object
ignored][Warning: Draw object ignored][Warning: Draw object
ignored][Warning: Draw object ignored][Warning: Draw object
ignored][Warning: Draw object ignored][Warning: Draw object
ignored][Warning: Draw object ignored][Warning: Draw object
ignored][Warning: Draw object ignored]Eines der traurigsten
Dinge im Leben ist,

\begin{figure}
\centering
 [Warning: Image ignored]
% Unhandled or unsupported graphics:
%\includegraphics[width=153.62mm,height=92.93mm]{FehlerartikelBuch-img/FehlerartikelBuch-img001.png}

\end{figure}
dass ein Mensch viele gute Taten tun muss,

um zu beweisen, dass er tüchtig ist,

aber nur einen Fehler zu begehen braucht,

um zu beweisen, dass er nicht taugt.

George Bernard Shaw




Fehler – Fehlermanagement - Sicherheitskultur




\section{Einleitung}

Auch wenn Organisationen eine Nullfehlerrate anstreben, so
geschehen trotzdem Fehler. Wie nämlich die Erfahrungen in
der Luftfahrtindustrie, der Atom- und chemischen Industrie
aber auch auf modernen Flugzeugträgern, wo man sich seit
Jahrzehnten damit beschäftigt, die Systeme sicherer zu
machen, zeigen, geschehen in komplexen Systemen, beinahe
zwanghaft Fehler, die jedoch im Einzelfall nicht unbedingt
einen Schaden verursachen müssen. Auch die Medizin ist so
ein komplexes System, in dem verschiedene Menschen in
verschiedenen Berufsgruppen mit hoch komplexen, technischen
Geräten und hoch potenten Medikamenten arbeiten.




Die Tatsache, dass auch in der Medizin Fehler geschehen und
gerade Ärzte und Ärztinnen nicht fehlerfrei arbeiten, wurde
zum ersten Mal 1999 durch das Buch \enquote*{To Err Is Human} des
Instituts of Medicine (Kohn/Corrigan et al. 1999) weltweit
bekannt und anerkannt. Die Autoren konnten zeigen, dass
circa vier von hundert stationären Patienten einen
Gesundheitsschaden während der Behandlung erleiden; weiters
konnten sie zeigen, dass von diesen Fehlern circa die
Hälfte vermeidbar gewesen wäre. Die Autoren schätzten die
Zahl der Todesfälle durch Fehler während der Behandlung mit
50.000 bis 100.000 pro Jahr höher ein als die Zahl der
Verkehrstoten in den USA in einem Jahr.




Ähnliche Ergebnisse erbrachte bereits 1995 die Untersuchung
\enquote*{the Quality in Australian Health Care Study}
(Wilson/Runciman et al., 1995, 458 – 471). 16,6\% der
Patienten waren von einem Adverse Event betroffen, der in
51\% vermeidbar gewesen wäre. Etwa drei Viertel der
Patienten hatten die Folgen dieses Ereignisses nach einem
Jahr überwunden, rund fünfzehn Prozent der Patienten trugen
bleibende Schäden davon, etwa jeder zwanzigste Patient
verstarb an dem Ereignis.




Diese Daten können wahrscheinlich problemlos auf die
Situation in Österreich übertragen werden. 2008 erschien in
\enquote*{Sozialversicherung online}, dem Organ des österreichischen
Hauptverbandes der Sozialversicherungsträger, eine Studie
zu diesem Thema. Die Autoren untersuchten 294.848 nach dem
LKF System (leistungsorientierte Krankenhausfinanzierung)
kodierte Krankengeschichten aus den Jahren 2001 – 2006 aus
Österreich. In 18.903 Fällen wurden Vergiftungen durch
Arzneimittel, Drogen und biologisch aktive Substanzen als
Diagnosen kodiert; in 105.908 Fällen wurden Komplikationen
bei chirurgischen Eingriffen und medizinischer Behandlung
kodiert. Die Autoren schätzen, dass 19.000 – 25.000 Fälle
oder 37\% - 51\% der Schäden vermeidbar gewesen wären.
Diese Fälle von unerwünschten Ereignissen hätten
Zusatzkosten von € 1,2 Milliarden verursacht
(Endel/Wilbacher 03/2008).




Gerade in Österreich ist der Umgang mit Fehlern
traditioneller Weise sehr stark von Verdrängen, Leugnen und
Vertuschen geprägt. Primär herrscht die Meinung vor, nur
sogenannt schlechte, unmotivierte und bös- beziehungsweise
unwillige MitarbeiterInnen machten Fehler. Besonders im
Gesundheitssystem, das aus Tradition stark hierarchisch
organisiert ist, sieht der Umgang mit Fehlern oft so aus,
dass zunächst ein Schuldiger gesucht und dieser dann
exemplarisch bestraft wird, was bei der betroffenen Person
beinahe automatisch zu defensivem Verhalten führt.




Riskmanagement beschäftigt sich daher oft auch
ausschließlich mit den haftungsrechtlichen Konsequenzen
eines Fehlers. Auf diese Art und Weise wird zwar die
Erwartungshaltung der Öffentlichkeit und der oder des von
den Folgen des Fehlers, dem Schaden, betroffenen PatientIn
befriedigt. Die Chance, einen Fehler auch als Information
über das System wahrzunehmen, wird damit jedoch vertan.




Ziel eines erfolgreichen Riskmanagements müsste es jedoch
sein, eine gute Fehlerkultur oder besser eine gute
Sicherheitskultur aufzubauen, denn Fehler sind in einem
System, in dem Menschen arbeiten, zwar unerwünscht und
unerfreulich, gleichzeitig aber auch unvermeidlich. Auch
wenn gerade in der Medizin die Folgen eines Fehlers immer
katastrophal bleiben werden (Tod oder schwere, eventuell
dauerhafte Schädigung des Patienten), so lässt sich doch
die Eintrittswahrscheinlichkeit für Fehler erhöhen. Um
dieses Ziel zu erreichen, wird die Information über latente
Fehler und Gefahren benötigt. Informationen über Fehler und
Beinahe – Fehler erlauben es so, gezielte und effektive
Maßnahmen zur Fehlervermeidung zu setzen. Diese Information
ist jedoch nur dann zu erhalten, wenn der / die einzelne
MitarbeiterIn nicht fürchten muss, für einen Fehler
abgestraft zu werden. Dann und nur dann werden Fehler nicht
länger vertuscht.




\section{Definitionen und Begriffe}

\subsection{Fehler}

Eine einheitliche Definition des Fehlers kann in der
Literatur nicht gefunden werden. Gründe dafür sind die
unterschiedliche Herangehensweise der jeweiligen Autoren,
die hohe Interdisziplinarität des Themas, die Abhängigkeit
von Aufgabencharakteristika, psychologischen
beziehungsweise umweltbedingten Faktoren (Löber 2012: 14).




Eine anerkannte Fehlerdefinition stammt vom englischen
Psychologen James Reason; für ihn sind mit Fehlern \enquote*{die
Ereignisse umfasst, bei denen eine geplante Abfolge
geistiger oder körperlicher Fähigkeiten nicht zum
beabsichtigen Resultat führt, sofern dieser Misserfolg
nicht fremden Einwirken zugeschrieben werden kann} (Reason
1994: 398).




Hofinger definiert Fehler als \enquote*{eine Abweichung von einem als
richtig angesehenen Verhalten oder von einem gewünschten
Handlungsziel, das der Handelnde eigentlich hätte ausführen
bzw. erreichen können} (Hofinger 2008: 37)




Gemeinsam ist diesen Definitionen, dass ein Fehler ein
geplantes Ziel voraussetzt, das durch eine aktive Handlung
nicht oder teilweise nicht erreicht wird; auch muss das
Ziel - zumindest theoretisch - erreichbar gewesen sein.




\subsection{Beinahe – Fehler, near - miss}

Dieser Begriff wurde aus der Luftfahrtindustrie übernommen.
Ein Beinahe – Fehler liegt dann vor, wenn durch glückliche
Umstände (z.B. rechtzeitiges Eingreifen oder Bemerken) ein
Fehler ohne Folgen, das heißt ohne Schaden, bleibt (Roth
2006: 15). Damit bleiben zum Schluss ein unerwünschtes
Ereignis und damit ein Schaden aus.




\subsection{unerwünschtes Ereignis, adverse event}

Ein unerwünschtes Ereignis ist die Folge eines Fehlers in
der medizinischen Versorgung, die dadurch zu einem Schaden
für den Patienten führt (Roth 2006: 15). Da auch
Nebenwirkungen von Medikamenten und Komplikationen von
chirurgischen Maßnahmen zu unerwünschten Ereignissen
führen, ist es sinnvoll, zwischen vermeidbaren (durch
Fehler verursachte) und unvermeidbaren unerwünschten
Ereignissen zu unterscheiden (z.B. Nebenwirkungen von
Arzneimitteln).




\subsection{technische Fehler}

Diese fassen alle Formen von Fehlfunktionen technischer
Geräte zusammen, die in Zusammenhang mit einem Fehler
verwendet wurden. Da die Folgen technischer Fehler meist
unmittelbar nach ihrem Auftreten bemerkt werden, sind diese
für die Zukunft verhältnismäßig leicht zu verhindern. Zu
ihrer Prävention existiert eine große Zahl an
Möglichkeiten, zum Beispiel die Fehler- Möglichkeiten- und
Einflussanalyse (FMEA) oder die Fehlerbaumanalyse (Roth
2006: 16).




\subsection{menschliche Fehler, der Human Factor}

Darunter werden alle diejenigen Fehler verstanden, die durch
Handlungen eines Menschen ausgelöst werden. Algedri und
Frieling (Algedri/Frieling 2001: 22) differenzieren
menschliche Handlungsfehler wie folgt

\begin{tabularx}{\linewidth}{XX}
\centering Handlung &
\centering\arraybslash Fehlerarten\\
Vorbereitung &
Informationsfehler

Wissensfehler

Wahrnehmungsfehler

Gedächtnisfehler\\
Durchführung &
Informationsfehler

Wissensfehler

Vertauschungsfehler

Hinzufügungs-, Auslassungs-, Positionierungsfehler

Reihenfolgefehler

Zeitpunktfehler

Mengenfehler\\
Kontrolle &
Informationsfehler

Wissensfehler

Urteilsfehler

Gedächtnisfehler\\
\end{tabularx}
Abbildung 2: Klassifikation menschlicher Handlungsfehler,
nach Algedri/Frieling 2001, eigene Darstellung




Untersuchungen aus Luftfahrt und Medizin konnten zeigen,
dass rund drei Viertel der Fehler durch den Human Factor
verursacht sind. Je nach Untersucher und Studie liegt die
Rate der Fehler, bei denen der Human Factor eine Rolle
spielt, zwischen 50\% und 83\% (Arnstein 1997).




Untersuchungen aus der Luftfahrtindustrie haben ergeben,
dass Fachwissen (technical skills) allein nicht ausreichen,
Non – Technical Skills sind genauso wichtig. Mit dem
Begriff Non – Technical Skills werden jene Fähigkeiten
bezeichnet, die über das konkrete Fachwissen hinausgehen,
Kommunikation, Fähigkeit zu Teamwork,
Situationsbewusstsein.\footnote{ Cf.:
http://de-wikipedia.org/wiki/Non-Technical\_Skills
(18.06.2013)} Fehler, die durch den Human Factor entstehen,
haben meist Defizite in den Non – Technical Skills zur
Ursache. Eine gute Möglichkeit für die / den einzelne/n
MitarbeiterIn den Human Factor an sich selbst zu erleben,
ist Crisis (Crew) Resource Management\footnote{ Crew
Resource Management wurde Ende der 1970 – er Jahre
entwickelt um die Sicherheit in der Luftfahrt zu erhöhen,
denn Untersuchungen hatten ergeben, dass menschliches
Versagen die Hauptursache für Flugunfälle war, die
Hauptprobleme waren Kommunikationsschwächen,
Kompetenzkonflikte und Entscheidungsschwäche. Seit gut zehn
Jahren wird Crew Resource Management als Crisis Resource
Management im medizinischen Training angewandt.\par Cf.:
http://de.wikipedia.org/wiki/Crew\_Resource\_Management
(18.06.2013)}, in dem die Non - Technical Skills zusätzlich
zu den Technical Skills beübt werden.




\subsection{Fehlermanagement}

Unter Fehlermanagement versteht man die systematische
Beschäftigung mit Fehlern und ihren Folgen, die in einem in
einem Mensch – Maschine – System auftreten können. Im
Rahmen des Fehlermanagements sollen die
Eintrittswahrscheinlichkeit für Fehler gesenkt und die
Auswirkungen eines Fehlers reduziert werden. Dazu sind die
systematische Fehlerbewertung sowie die aktive Suche nach
potentiellen Fehlern und Gefahren notwendig\footnote{ Cf.:
http://de.wikipedia.org/wiki/Fehlermanagement,
(01.03.2013)}. Ähnlich wie im Qualitätsmanagement wird im
Fehlermanagement nach dem PDCA – Zyklus (Deming –
Zyklus)\footnote{ Cf.:
http://de.wikipedia.org/wiki/Demingkreis, (01.03.2013)}
vorgegangen (Roth 2006: 40)

Plan: Fehler kontinuierlich verringern durch
Fehlermanagement

Do: Fehler identifizieren und melden

Check: Fehlerevaluation durch Fehleranalyse

Act: Fehlerreduktion durch Verhaltenstraining und
Verbesserung

Im darauf folgenden Durchgang wird anstatt \enquote*{Plan} eine
Reevaluation vorgenommen, deren Ergebnisse dann in den
neuen Zyklus einfließen.







\begin{figure}


Do

Check

Act


Plan

\caption{Abbildung 3: Demingzyklus, eigene Darstellung}
\end{figure}
























Um ein Fehlermanagementsystem wirksam werden zu lassen,
stellt Leape folgende Anforderungen an ein
Fehlermeldesystem (Leape 2002: 1633 – 1638, eigene
Übertragung) 



\begin{table}
\begin{tabularx}{\linewidth}{XX}
\centering Eigenschaft &
\centering Erklärung
\\
sanktionsfrei &
Keine Bestrafung des/r Meldenden
\\
vertraulich &
Die Identität des/r Meldenden und aller Beteiligten wird
nicht an dritte weitergegeben
\\
eigenständig &
Die Meldung ist unabhängig von jeder Autorität, die den
Meldenden bestrafen könnte
\\
Analyse durch Experten &
Die Meldung wird von einem Expertenteam analysiert, das die
klinischen Zusammenhänge versteht und Erfahrung hat, die
systemischen Ursachen zu erkennen
\\
zeitnah &
Berichte werden schnell analysiert, die Ergebnisse werden
schnell an die Betroffenen weitergegeben, besonders, wenn
Gefahr im Verzug ist
\\
systemorientiert &
Die Empfehlungen beziehen sich auf Veränderungen von
Systemen, Prozessen und Verfahrensanweisungen
\\
verantwortlich &
Der Empfänger der Meldungen hat sowohl die Macht als auch
die Bereitschaft, Vorschläge umzusetzen
\\
\end{tabularx}
\caption{}
\end{table}



\subsubsection{Aktive versus latente Fehler}

Aktive Fehler sind solche sicherheitsgefährdende Handlungen,
deren Folgen meistens unmittelbar wahrgenommen werden. Bei
latenten (passiven) Fehlern kommt es in der Regel nicht
unmittelbar zu einem Schaden; es handelt sich um Fehler,
Irrtümer oder unsichere Handlungen, die oft lange bestehen
können, ohne deswegen gleich einen Schaden zu verursachen.
Erst das Zusammentreffen von mehreren latenten und / oder
aktiven Fehlern, beziehungsweise das Versagen von
Sicherheitsbarrieren führt zu einem konkreten Schaden. Eine
der Aufgaben des Fehlermanagements ist es unter anderen,
nach diesen latenten Fehlern zu fahnden (Roth, 2006: 19).




Das Verhältnis von latenten Fehlern und großem,
katastrophalen Schaden zeigt Heinrichs Gesetz auf. Eine
Untersuchung des Ingenieurs Heinrich 1941 konnte zeigen,
dass sich bei 3846 Patienten etwa 300 Fehler (Beinahe -
Unfälle) ohne oder nur mit sehr geringem Schaden ereignen;
29 Ereignisse waren mittelschwer und verursachten einen
entsprechenden Schaden; nur in einem Fall trat das
katastrophale Schadensereignis ein. Heinrichs Gesetz zeigt
also sehr deutlich, wie wichtig es ist, nach den 300
Fehlern ohne wesentlichen Schaden zu fahnden\footnote{ Cf.:
http://duepublico.uni-duisburg-essen.de/servlets/DerivateServlet/Derivate-26767/Dissertation\%20Risikomanagement\%20im\%20Gesundheitswesen.pdf
(01.03.2013)}, um den einen katastrophalen Schaden zu
vermeiden. Abbildung 4 veranschaulicht Heinrichs Gesetz.


\begin{figure}
Abbildung 4: 
\caption{Heinrichs Gesetz, eigene Darstellung}
\end{figure}








\subsubsection{Personenorientierter versus systemorientierter
Ansatz}

Fehler können aus zwei Perspektiven
analysiert werden, dem personenorientieren Ansatz
(Leape/Woods et al., 1998: 1444 – 1447; Reason: 2000, 768 -
770) und dem systemorientierten Ansatz (Bates/Cullen et
al., 1995: 29 – 34; Reason, 2000: 768 – 770). Jeder Ansatz
hat seine eigene Vorstellung davon, wie Fehler entstehen.
Auch werden in jedem Ansatz andere Strategien zur
Fehlervermeidung vorgeschlagen.

Der
personenorientierte Ansatz ist in gewisser Weise die
traditionelle Sicht des Fehlers, die Ereignisse
ausschließlich vom aktiven Fehler her analysiert. So wird
der / die einzelne MitarbeiterIn, bei dem / der ein Fehler
– \enquote*{so etwas} – passiert ist, als alleinige/r Schuldige
darstellt. Der personenorientierte Ansatz sucht die
Fehlerursachen in Schlamperei, Nachlässigkeit,
Unachtsamkeit, Unwissen, mangelndem Engagement und
ähnlichem der handelnden Personen. Die allseits anerkannte
These dabei ist, sogenannt gute MitarbeiterInnen machten
keine Fehler. Deshalb zielen auch alle Maßnahmen zur
Fehlervermeidung auf den einzelnen Mitarbeiter, im
Einzelnen disziplinäre Maßnahmen, Nachschulungen,
Klagsdrohungen, in Schlagworten zusammengefasst als name
it, blame it, shame it; die Culture of Blame.

Gerade im sehr hierarchisch
strukturierten Gesundheitswesen ist der traditionelle,
personenorientierte Ansatz weit verbreitet. Es ist eben
befriedigender (\enquote*{satisfying}, Reason 2000: 768 - 770),
Menschen zu beschuldigen und zu tadeln, als sich mit den
Strukturen der Organisation zu beschäftigen (Reason, 2000:
768 – 770). Eine solche drohende Haltung schüchtert die
MitarbeiterInnen ein und führt letztendlich dazu, dass sich
die MitarbeiterInnen insgesamt ausschließlich defensiv
verhalten. Fehler werden geleugnet und versteckt; von Eiff
spricht in diesem Zusammenhang von der \enquote*{verborgenen Fabrik}
(von Eiff/Stachel, 2007: 172). Eine systematische
Fehlersuche wird damit allerdings unmöglich gemacht oder
zumindest erschwert.

Im systemorientierten Ansatz hingegen
wird davon ausgegangen, dass Menschen von Natur aus Fehler
machen und Fehler nur selten oder beinahe nie der
Böswilligkeit oder dem Nichtwissen des/r MitarbeiterIn
entspringen. Da in diesem Ansatz der Human Factor anerkannt
wird, werden selbst im besten System, solange Menschen
darin arbeiten, Fehler geschehen.

Reason beschreibt Fehler als Folgen der
Systemstruktur und nicht als Ursache für Ereignisse. Er
geht davon aus, dass die menschliche Natur, die Human
Factors, eben nicht verändert werden können. Fehler werden
nach Reason nicht als Folgen von Nachlässigkeit oder
menschlichen Unzulänglichkeiten interpretiert. Aus diesem
Grund sei es notwendig, die Bedingungen unter denen
Menschen arbeiten, zu ändern (\enquote*{we cannot change human
conditions, but we can change the conditions under which
humans work}. Reason, 2000: 768 - 770). Ziele sind die
Entwicklung eines Systems, das das Entstehen von Fehlern
erschwert, und die Konstruktion von
Sicherheitsvorkehrungen, die die Auswirkungen von trotzdem
eingetretenen Fehlern vermindern (Reason, 2000: 768 – 770)



Im systemorientierten Ansatz werden die Barrieren in der
Organisation, die Fehler verhindern sollen, ins Zentrum der
Analyse gestellt; die Frage lautet nicht, wem ist so etwas
passiert, sondern was ist passiert und warum waren die
Sicherheitsmaßnahmen nicht ausreichend. Reason
veranschaulicht diesen Ansatz mit seinem \enquote*{Swiss Cheese
Model} (Abbildung 5).





\begin{figure}
  [Warning: Image ignored]
% Unhandled or unsupported graphics:
%\includegraphics[width=153.63mm,height=92.92mm]{FehlerartikelBuch-img/FehlerartikelBuch-img002.png}
 

\centering
\begin{minipage}{39.99mm}
Unerwünschtes Ereignis, Schaden
\end{minipage}
\centering
\begin{minipage}{45mm}
Lokale Auslöser, untypische Bedingungen
\end{minipage}
\centering
\begin{minipage}{39.99mm}
Latentes Versagen
\end{minipage}
\centering
\begin{minipage}{29mm}
Gefahr
\end{minipage}
\centering
\begin{minipage}{39.99mm}
Unsichere Handlungen
\end{minipage}
\centering
\begin{minipage}{39.99mm}
Psychologische Vorläufer
\end{minipage}
\centering
\begin{minipage}{39.99mm}
Sicherheitsbarrieren
\end{minipage}
\caption{Abbildung 5: Swiss Cheese Model nach Reason, eigene
Darstellung}
\end{figure}









Idealer Weise wären die Barrieren, die die Organisation
errichtet hat, um Fehler zu vermeiden, völlig intakt; in
Wirklichkeit haben sie jedoch Löcher, die je nach
Einzelfall abhängig von unvorhersehbaren Einflüssen sowohl
innerer als auch äußerer Faktoren, offen oder geschlossen
sind, beziehungsweise ihre Position verändern. Im Regelfall
führt das Versagen einer einzelnen Barriere noch nicht zu
einem konkreten Schaden. Das Entstehen eines Schadens, eine
Katastrophe, setzt fast immer die Kombination eines aktiven
Fehlers mit latenten Konditionen (latent conditions)
voraus. Erst wenn die Löcher, entstanden aus aktivem Fehler
und latenten Konditionen, in einer fatalen Ebene liegen,
entsteht ein Unfall. Den Zusammenhang zwischen aktivem
Fehler und latenten Konditionen, den Fehlerebenen und
Fehlerursachen, verdeutlich Abbildung 6.




\clearpage



\begin{figure}

Fehlerebenen und Fehlerursachen


Patienten-schaden
Sicherheitsbarriere
Sicherheitsbarriere
Individual-faktoren

Physiologische/biologische Faktoren

Wissens- \& fertigkeits-bezogene Faktoren

Kognitive \& motivationale Faktoren

Aktive Fehler

Latente Fehler

Latente Fehler

Latente Konditionen

Management-faktoren













Organisations-prozesse

Fehlerkultur

Trigger-faktoren

Arbeitsum-gebung

Teamfaktoren

Arbeitsbe-zogene Faktoren

Patient

Technik

Systemische Kontext-faktoren im Krankenhaus

Dynamik

Intransparenz

Situative Vernetztheit

Unsicherheit

Temporo-spatiale Konfiguration

\caption{Abbildung 6:: Fehlerebenen und Fehlerursachen, modifiziert
nach Vincent C. et al., British Medical Journal 2000 in
Löber 2012}
\end{figure}


Das Swiss Cheese Model von Reason erlaubt so, das Entstehen
von fatalen Ereignissen als Zusammenwirken von
Entscheidungen des höchsten Managements mit Handlungen
des/r einzelnen MitarbeiterIn zu verstehen (Löber 2012: 46)




\subsection{Fehlerkultur, Sicherheitskultur}

Der Begriff Fehlerkultur wurde im Wesentlichen nach dem Jahr
2000 im deutschsprachigen Raum entwickelt; in der
englischsprachigen Literatur wird meist der Begriff Safety
Culture, Sicherheitskultur, verwendet. Fehlerkultur ist ein
Teil der Unternehmenskultur, also der Art und Weise, wie
MitarbeiterInnen und Führungskräfte miteinander
kommunizieren. Die Unternehmenskultur definiert sich durch
Symbole, Helden, Rituale und Werte (von Eiff/Stachel, 2007:
18ff.).




Mit Fehlerkultur wird der geänderte Umgang mit Fehlern
beschrieben. Der Begriff bezeichnet die Art und Weise, wie
ein System mit Fehlern, Fehlerrisiken und Fehlerfolgen
umgeht (Schüttelkopf, 2008: 233). Die oberflächliche
Sichtweise der individuellen Schuldzuweisung (Culture of
Blame) soll zu Gunsten einer systemanalytischen, proaktiven
Sicherheitskultur (Safety Culture) verlassen werden, Ziel
ist der vorurteilsfreie Umgang mit Fehlern
(Holzer/Thomeczek et al. 2005: 166). Es geht darum, ob
Fehler und ihre Folgen als Chancen zum Lernen wahrgenommen
oder ausschließlich negativ bewertet werden (Lindemann
2005: 25). Fehler sollen in ihrem Entstehungskontext
wahrgenommen und beurteilt werden (Abed – Navandi 2007:
103).




Mit Fehlerkultur wird so etwas wie ein gewünschter Zustand
beschrieben, der einer konstruktiven Fehlerkultur ähnlich
ist. Es geht um den Umgang mit und die Beurteilung von
Fehlern genauso wie um das Aufspüren potentieller Fehler.
Was allerdings in all diesen Definitionen fehlt, ist der
Bezug zur Unternehmenskultur (Löber 2012, 191ff.). Daher
soll an dieser Stelle die Definition für Sicherheitskultur
der britischen \enquote*{Health and Safety Commission} ergänzt
werden. \enquote*{Fehlerkultur als Teilkonstrukt der
Unternehmenskultur ist das Produkt individueller und
kollektiver Werte, Einstellungen, Kompetenzen und
Verhaltensmuster, die das Ausmaß, die Art und die Tiefe der
organisationalen Auseinandersetzung mit innerbetrieblichen
Fehlern bestimmen} (Health and Safety Commission, 1993:
23). Manstein fasst es so zusammen: Sicherheitskultur ist
\enquote*{the way we do things around here} (Manstein 1999).




Fehlerkultur ist also Teil der Unternehmenskultur und wird
von dieser beeinflusst. Im Besonderen wird die Fehlerkultur
von den Werten der Organisation beeinflusst. In dem Maß als
Kultur aktiv gestaltbar ist, ist auch Fehlerkultur
veränderbar (Löber 2012, 194f).




Reason legt ein Modell der Sicherheitskultur zur Entwicklung
eines betrieblichen Fehlerkulturmodells vor, das durch vier
Elemente charakterisiert ist (Reason 1997, 197f.)

\begin{itemize}
  \item Reporting Culture
  \item Just Culture
  \item Flexible Culture
  \item Learning Culture
\end{itemize}
Eine Sicherheitskultur kann nur erreicht werden, wenn alle
MitarbeiterInnen gleichen Informationsstand über Beinahe –
Unfälle und Fehler haben oder wenn solche Informationen
abrufbar sind; Reason verwendet den Begriff Reporting
Culture als Synonym für Reporting System.




Just Culture beschäftigt sich mit dem Umgang mit geschehenen
Fehlern und beinahe – Fehlern. Es geht um die Motivation
der MitarbeiterInnen, im Rahmen der Reporting Culture
solche Ereignisse zu melden. Just Culture ist für das
Funktionieren einer Sicherheitskultur zentral. Persönliche
Haftung und Verantwortlichkeit auf der einen Seite müssen
mit der gewünschten Sicherheit auf der anderen Seite
ausbalanciert werden.




Komplexe Systeme zeichnen sich durch schnell wechselnde
Anforderungen aus. Daher ist Flexible Culture für eine
Sicherheitskultur wichtig, um auf diese wechselnden
Bedingungen rasch und passend reagieren zu können.
Hochrisikoorganisationen (wie sie in Kapitel 2.4
beschrieben werden) haben hohe Kompetenz auf dem Gebiet der
Flexible Culture.




Learning Culture bedeutet den konstruktiven Umgang mit
Fehlern. Fehler und die Analyse von Fehlern werden in einer
guten Sicherheitskultur als Chance wahrgenommen,
Informationen über die Stärken und Schwächen der
Organisation zu erhalten; das Stichwort heißt, Lernen aus
Fehlern.




Wie ein modernes Konzept der Fehlerkultur (Löber 2012: 219)
aussieht, zeigt Abbildung 7.









\begin{figure}

\centering
\begin{minipage}{44.49mm}
Kultureller Rahmen
\end{minipage}

\centering
\begin{minipage}{50.78mm}
Fehlerkultur
\end{minipage}

\centering
\begin{minipage}{26mm}
Kommuni-kation
\end{minipage}

\centering
\begin{minipage}{26mm}
Positive Emotionen
\end{minipage}


\centering
\begin{minipage}{26mm}
Gerechtig-keit
\end{minipage}

\centering
\begin{minipage}{26mm}
Lernen
\end{minipage}

\centering
\begin{minipage}{26mm}
Vertrauen
\end{minipage}

\centering
\begin{minipage}{26mm}
Flexibilität
\end{minipage}


\caption{Abbildung 7: Modell der Fehlerkultur, Löber 2012}
\end{figure}


Hoffmann und Rohe haben 2010 ein fünfstufiges System zur
Beurteilung der Reife der Sicherheitskultur einer konkreten
Organisation oder Organisationseinheit vorgelegt, die fünf
Reifegrade der Sicherheitskultur (Hoffmann/Rohe, 2010: 94).

\begin{itemize}
  \item Reifegrad 1, ablehnende Sicherheitskultur
  \item Reifegrad 2, reaktive Sicherheitskultur
  \item Reifegrad 3, vorschriftsmäßige Sicherheitskultur
  \item Reifegrad 4, initiative Sicherheitskultur,
  \item Reifegrad 5, zukunftsweisende Sicherheitskultur
\end{itemize}



In einer ablehnenden Sicherheitskultur werden bei kritischen
Ereignissen vor allem Schuldige gesucht, um diese zu
bestrafen. Patientensicherheit wird nur zögerlich
betrachtet. Ein Lernen aus Fehlern findet kaum statt. In
der reaktiven Sicherheitskultur wird nur dann gehandelt,
wenn ein Schaden geschehen ist. Erfüllen von externen
Vorgaben wie zum Beispiel die Einführung von Qualitäts- und
Riskmanagementsystemen ist das Kennzeichen der
vorschriftsmäßigen Sicherheitskultur. Gibt es Maßnahmen zur
Verbesserung der Patientensicherheit ohne konkreten
Anlassfall, so hat die Organisation den Reifegrad der
initiativen Sicherheitskultur erreicht. Wenn schließlich
Patientensicherheit selbstverständliches Anliegen aller
MitarbeiterInnen und routinierter Teil ihres täglichen
Handelns ist, so liegt eine zukunftsweisende
Sicherheitskultur vor.




\subsection{High Reliability Organisation, die
Hochrisikoorganisation, HRO}

High Reliability Organisations oder Hochrisikoorganisationen
sind solche Organisationen, in denen das Auftreten von
Schäden erfolgreich vermieden wird, oder die Fehlerrate
niedriger als erwartet ist, obwohl permanent in einem hoch
komplexen und gefährlichen Umfeld gearbeitet wird. Seit
etwa 1990 beschäftigt sich die Wissenschaft mit diesen
Organisationen.\footnote{ Cf.:
http://en.wikipedia.org/wiki/High\_reliability\_organization
(02.03.2013)} Im Mittelpunkt des Forschungsinteresses
stehen die sich entwickelnde wechselseitige Beziehung
zwischen Verhalten und Struktur der Organisation. Beispiele
für solche Hochrisikoorganisationen sind die Atomindustrie,
die chemische Industrie, die Luftfahrtindustrie,
Flugzeugträger genauso wie Rettungsorganisationen und
Krankenhäuser (Mistele 2005: 4)




Für die Zuverlässigkeit einer Organisation ist die
Unternehmenskultur von besonderer Bedeutung. Zentral für
die Hochrisikoorganisation ist der Begriff der
\enquote*{mindfulness}, Achtsamkeit. Diese ist gekennzeichnet durch

\begin{itemize}
  \item Sensibilität für betriebliche Abläufe
  \item Konzentration auf Fehler
  \item Abneigung gegen vereinfachende Interpretationen
  \item Streben nach Flexibilität
  \item Respekt vor fachlichem Können und Wissen.
\end{itemize}



Mindfulness hilft Fehlverläufe frühzeitig – bevor noch
größerer Schaden eingetreten ist – zu erkennen.
(Sensibilität für Abläufe, Aufmerksamkeit für Fehler und
Gefahren, Anerkennen der komplexen Vorgänge). Gleichzeitig
übt Achtsamkeit den Umgang mit Unerwartetem (Streben nach
Flexibilität und Respekt vor Fachwissen) (Mistele 2005: 13)







\section{Literatur}

\begin{description}
\item[AbedNavandi-NullFehler] Abed – Navandi, M.: Null – Fehler ja, Null –
Fehlerkultur nein, Fehlerkultur, eine Notwendigkeit ohne
Alternative, in: Wagner, R. (Hrsg.), Near Miss,
systematischer Umgang mit Beinahe – Unfällen, London 2007
\item[] Algedri, J./Frieling, E.: Human FMEA: menschliche
Handlungsfehler erkennen und vermeiden, München/Wien 2001
\item[] Arnstein, F.: British Journal of Anesthesia 1997, 79:
645 ff
\item[] Bates, D./Cullen, D. et al.: Incidence of Adverse Drug
Events and Potential Adverse Drug Events. Implications for
Prevention. ADE Prevention Study Group, Journal of the
American Medical Association 1995; 274 (1): 29 – 34
\item[] Endel G./Wilbacher I.: Patientensicherheit:
Unerwünschte Ereignisse, Sozialversicherung online, 03/2008
\item[] Health and Safety Commission, Third Report: Organizing
for Safety, ACSNI Study Group on Human Factors – HMSO,
London 1993, zitiert nach Löber 2012: 193
\item[] Hoffmann, B./Rohe, J.: Patientensicherheit und
Fehlermanagement, Ursachen unerwünschter Ereignisse und
Maßnahmen zu ihrer Vermeidung, Übersichtsarbeit, Deutsches
Ärzteblatt 2010, 107 (6): 92 - 100
\item[] Hofinger, G.: Fehler und Unfälle, in Badke – Schaub,
P./Hofinger, G. et al. (Hrsg.): Human Factors: Psychologie
sicheren Handelns in Risikobranchen, Heidelberg 2008
\item[] Holzer, E./Thomeczek, C. et al.: Patientensicherheit,
Leitfaden für den Umgang mit Risiken im Gesundheitswesen,
Wien 2005
\item[] Kohn, L./Corrigan, J. et al. (eds.): To Err Is Human,
Building a Safer Health System. Committee on Quality of
Health Care in America. Washington DC, Institute of
Medicine 1999
\item[] Leape, L./Woods, D. et al.: Promoting Patient Safety
by Preventing Medical Error, Journal of the American
Medical Association, 1998; 280 (16): 1444 – 1447
\item[] Leape, L.: Reporting of Adverse Events, New England
Journal of Medicine, 2002; 347 (20): 1633 – 1638
\item[] Lindemann, U.: Methodische Entwicklung technischer
Produkte: Methoden flexibel und situationsgerecht anwenden,
Berlin 2005
\item[] Löber, N.: Fehler und Fehlerkultur im Krankenhaus,
eine theoretisch – konzeptionelle Betrachtung, Wiesbaden
2012
\item[] Manstein, Z.: Organizational Culture and Safety
Performance. Does Your Corporate Culture Promote Safety?
How Can You Make Safety Part of the Corporate Culture, EHS
Today 1999, http://ehstoday.com/news/ehs\_imp\_32830
(05.03.2012)
\item[] Mistele, P.: die Relevanz der High Reliability Theory
für Hochleistungssysteme, Diskussionspapier, Chemnitz 2005,
http://www.tu-chemnitz.de/wirtschaft/bwl6/publikationen/publikation\_download.php?NR=175
(02.03.2013)
\item[] Reason, J.: Menschliches Versagen: Psychologische
Risikofaktoren und moderne Technologien, Heidelberg 1994
\item[] Reason, J.: Managing the Risks of Organizational
Accidents, Aldershot 1997
\item[] Reason, J.: Human Error, Models and Management,
British Medical Journal, 2000, 320 (7237): 768 – 770
\item[] Roth, A.: Fehlermanagement im Krankenhaus, Konzept zur
Implementierung eines Fehlerverständnisses, Saarbrücken
2006
\item[] Schüttelkopf, E.: Erfolgsstrategie Fehlerkultur, in
Ebner, Heimerl et al. (Hrsg.), Fehler – Lernen –
Unternehmen, wie sie die Fehlerkultur und Lernreife ihrer
Organisation wahrnehmen und gestalten, Frankfurt am Main
2008
\item[] Wilson, R./Runciman, W. et al.: the Quality in
Australian Health Care Study, the Medical Journal of
Australia, 1995 (163): 458 – 471
\item[] Vincent, C./Tayler – Adams, S. et al.: Framework for
Analyzing Risk and Safety in Clinical Medicine, British
Medical Journal 1998, 316: 1154 - 1157
\item[] von Eiff, W./Stachel, K. (Hrsg.): Unternehmenskultur
im Krankenhaus, Band 1 der Reihe \enquote*{leistungsorientierte
Führung und Organisation im Gesundheitswesen}, Gütersloh
2007
\end{description}


\ChapterEnd